%! Simplifying Algebraic Expressions and Exponent Rules — Unit 1, Lesson 1.1 — Student Packet Cover Sheet
\documentclass[10pt]{article}
\usepackage{atda-article}
\usepackage{atda-boxes}
\usepackage{ltablex}
\keepXColumns

\begin{document}

% Full-bleed royal banner
\begin{tikzpicture}[remember picture, overlay]
  \fill[royal] (current page.north west)
    rectangle ([yshift=-0.9in]current page.north east);
\end{tikzpicture}
\vspace{-0.8in}

\begin{center}
  {\color{white}\textbf{\LARGE Algebra, Trigonometry, and Data Analysis}}\\[4pt]
  {\color{mistmid}\large\textbf{Unit 1: Algebraic Expressions and Factoring}}\\[6pt]
  {\color{white}\textbf{\large Lesson 1.1 \quad Simplifying Algebraic Expressions and Exponent Rules}}
\end{center}

\vspace{0.22in}
\namedateperiod
\vspace{0.18in}

\begin{learningtargetbox}
\small
\begin{enumerate}[leftmargin=1.5em, itemsep=5pt]
  \item \ldots{} simplify a monomial expression in one or two variables using the
        \textbf{exponent rules}: product, quotient, power of a power, power of a product, and
        the zero and negative exponents.
  \item \ldots{} add and subtract polynomials by combining like terms, distributing the
        subtraction to \emph{every} term of the second polynomial.
  \item \ldots{} multiply polynomials --- monomial by polynomial, binomial by binomial in one
        or two variables, and binomial by trinomial --- and write the result in standard form.
  \item \ldots{} build a polynomial expression from a situation and use it to answer a
        question about that situation, then say what the answer means.
\end{enumerate}
\end{learningtargetbox}

\vspace{0.14in}

\begin{tocbox}
\small
\renewcommand{\arraystretch}{1.5}
\begin{tabularx}{\linewidth}{@{} c l X r @{}}
\rowcolor{mist}
\textbf{\#} & \textbf{Component} & \textbf{Description} & \textbf{Score} \\
\midrule
1 & Warm-Up        & Exponents as repeated multiplication; like terms; distributing
                                                                    & \blank{1.2cm} \\
2 & Guided Notes   & The exponent rules, sums and differences, and products of polynomials
                                                                    & \blank{1.2cm} \\
3 & Group Activity & The school store: a revenue, cost, and profit model & \blank{1.2cm} \\
4 & Exit Ticket    & A quotient, a difference, and one error to diagnose & \blank{1.2cm} \\
5 & Homework       & Mixed fluency, a concession-stand profit model, and scaling
                                                                    & \blank{1.2cm} \\
\midrule
  & \multicolumn{2}{r}{\textbf{Total}} & \blank{1.2cm} \\
\end{tabularx}
\end{tocbox}

\vspace{0.14in}

\begin{remindbox}
\small
This lesson covers \texttt{A2.EO.3a} --- determining sums, differences, and products of
polynomials in one and two variables. The exponent rules are the machinery underneath: they are
Algebra 1 and 2 review here, and you will spend them again in Lesson 1.2 (rational exponents)
and in Unit 3, where $\log(ab)=\log a+\log b$ is the product rule written backwards.
\end{remindbox}

\end{document}
